\section{A foreign subsidiary from acquisition to disposal}
\label{sec:group-fx-disposal}

The preceding examples establish acquisition, translation and disposal
separately. Their interaction matters because a foreign subsidiary carries
more than translated net assets into a sale: goodwill and accumulated exchange
differences also have owners and dates. A sale can leave total equity correct
while putting the wrong amount in profit. We follow one subsidiary through
these steps using the selected historical IFRS 3, IFRS 10, IAS 21 and IAS 7
provisions in the retained EU text \citep{eu20231803}.

A parent starts with USD 300 of cash and equity. USD is its functional and
presentation currency; the subsidiary's functional currency is EUR. At
acquisition the parent pays EUR 96 for 80\% of equal ordinary interests, when
each euro buys USD 1.20. It obtains substantive power over the relevant
activities, exposure to variable returns and the ability to affect them. The
other owners have no substantive rights that defeat that control. Neither
investment-entity treatment nor a common-control transaction applies.

The acquired business already produces outputs and includes a skilled
organised workforce performing a process critical to continuing them, as
required by the assumed IFRS 3.B12C(a) facts. The optional concentration test
is not elected. Identifiable fair values are EUR 20 of cash, inventory of
140 and a payable of 60. There is no other separately identifiable asset,
tax difference, previous holding, contingent consideration or acquisition cost
in this synthetic case. The workforce itself is subsumed into goodwill under
B37; calling a residual goodwill would not justify omitting an identifiable
customer relationship or other asset in a real acquisition.

Identifiable net assets are \(20+140-60=100\) euros. The parent chooses
proportionate-share measurement for the non-controlling interest (NCI) under
IFRS 3.19. NCI is therefore EUR 20 and goodwill under paragraph 32 is
\(96+20-100=16\) euros. This recognised goodwill belongs to the acquired
parent interest; no NCI goodwill has been recognised. The parent's separate
entry debits investment and credits cash by USD 115.2. The consolidation
adjustment debits acquisition equity 120 and goodwill 19.2, and credits the
investment 115.2 and NCI 24. Immediately after acquisition, group cash is
208.8, inventory 168, goodwill 19.2 and the payable 72. Net assets of 324
equal parent equity 300 plus NCI 24.

The subsidiary then sells all the inventory for EUR 160, receives cash, and
pays its EUR 60 payable. Both events occur at USD 1.25 per euro. The specified
event interval has no further expense, impairment or distribution. In the
subsidiary's books, revenue 160 less cost of sales 140 gives profit 20; cash
closes at \(20+160-60=120\) euros. Translated revenue is 200, cost of sales
175 and profit 25. Parent and NCI shares of profit are 20 and 5. The
translation worksheet translates the opening balances to the event rate
before recording these movements; it does not charge cost 175 against
unadjusted acquisition inventory 168 and leave a temporary negative stock.
These are presentation adjustments, not a remeasurement of historical-cost
inventory in the subsidiary's own functional currency.

At reporting, and at an immediate disposal considered below, the rate is
USD 1.30 per euro. Subsidiary cash translates to 156. IAS 21.47 also treats
the acquisition goodwill as an asset of the foreign operation: its EUR 16
becomes USD 20.8. To explain the currency reserve, translate the opening net
assets and the subsequently earned profit from their respective rates:
\[
  \underbrace{100(1.30-1.20)}_{\text{opening net assets}}
  +\underbrace{20(1.30-1.25)}_{\text{later profit}}=11.
\]
Goodwill adds \(16(1.30-1.20)=1.6\). The total translation gain in other
comprehensive income (OCI) is 12.6, but it is not divided mechanically 80/20.
Only the net-asset gain is shared: parent OCI is \(0.8(11)+1.6=10.4\),
and NCI OCI is \(0.2(11)=2.2\). This follows from the chosen proportionate
NCI measurement and the separate attribution required by IAS 21.41 and
IFRS 10.B94. NCI closes at \(24+5+2.2=31.2\), equal to 20\% of the
subsidiary's identifiable net assets of 156. Parent equity is
\(300+20+10.4=330.4\). Group cash 340.8 plus goodwill 20.8 equals their
combined equity of 361.6.

Cash has its own translation bridge. At acquisition, payment of 115.2 less
cash acquired of 24 gives investing outflow 91.2. The subsequent operating
receipt of 200 less payment of 75 gives operating inflow 125. Exchange
changes increase the translated value of the acquired cash by
\(20(1.30-1.20)=2\) and of the later EUR 100 net receipt by
\(100(1.30-1.25)=5\). IAS 7.25--28 reports that total exchange effect of
7 separately from transaction cash flows. Thus
\(300-91.2+125+7=340.8\). The cash exchange effect is 7; the total
translation OCI is 12.6. They measure different populations.

\subsection{Losing control removes the whole subsidiary}
\label{sec:foreign-disposal-reserves}

Suppose an unanticipated transaction immediately after reporting transfers
control to a buyer. No earlier held-for-sale condition is assumed in this
event slice; a real planned disposal needs that assessment. Compare two
alternatives: sell the whole 80\% interest for USD 169, or sell 60\% for
126.75 and retain 20\% with independently admitted fair value 42.25. Both
transactions remove the seller's substantive power. The equal price per
share is an illustrative valuation assumption, not an accounting requirement.
It makes proceeds plus retained value equal to 169 in each alternative.

IFRS 10.B98 removes all subsidiary assets, liabilities and NCI when control
is lost, and measures any retained investment at fair value. If \(C\) is
cash consideration and \(V\) the retained fair value, the gain before the
currency reserve is reclassified is
\[
  C+V+\underbrace{31.2}_{\text{NCI removed}}
      -\underbrace{156}_{\text{subsidiary cash removed}}
      -\underbrace{20.8}_{\text{goodwill removed}}=23.4.
\]
For the full sale, debit cash 169 and NCI 31.2, credit subsidiary cash 156
and goodwill 20.8, and credit gain 23.4. In the retained-interest alternative,
replace the debit to cash 169 with cash 126.75 and retained investment 42.25.
The subsidiary still holds its own bank account; removal of cash 156 is a
change in the consolidated reporting population, not a payment by the parent.

IAS 21.48--48B and IFRS 10.B99 then reclassify the parent's entire accumulated
exchange reserve of 10.4 into profit: debit that reserve and credit disposal
gain. Total disposal gain is \(23.4+10.4=33.8\). NCI's reserve of 2.2
is already within the NCI balance of 31.2 removed above; it is derecognised
without being recycled to profit. Retaining 20\% does not justify recycling
only the fraction corresponding to the shares sold. Paragraph 48A treats a
loss of control as a disposal even when a minority interest remains.

The two alternatives therefore have equal total ending equity but different
cash and investment positions:

\begin{center}
\begin{tabular}{lrr}
\toprule
USD amount & Sell all 80\% & Sell 60\%, retain 20\% \\
\midrule
Cash consideration & 169 & 126.75 \\
Retained investment & 0 & 42.25 \\
Disposal gain including recycling & 33.8 & 33.8 \\
Net disposal cash flow & 13 & $-29.25$ \\
Closing group cash & 353.8 & 311.55 \\
Closing equity & 353.8 & 353.8 \\
\bottomrule
\end{tabular}
\end{center}

IAS 7.39--42 reports acquisition and loss-of-control cash separately, each
net of the whole subsidiary's cash. Net disposal cash is \(169-156=13\)
for the full sale, but \(126.75-156=-29.25\) when 20\% is retained.
A positive parent receipt can therefore accompany a negative consolidated
disposal cash flow. The retained investment's 42.25 is not cash. The two
cash bridges are \(300-91.2+125+7+13=353.8\) and
\(300-91.2+125+7-29.25=311.55\).

For the whole interval, operating profit 25 plus disposal gain 33.8 gives
profit 58.8, attributed 53.8 to the parent and 5 to NCI. Total OCI is
\(12.6-10.4=2.2\), all attributed to NCI; period OCI does not vanish
merely because closing NCI is zero. The equity roll-forward is
\[
  300+\underbrace{24}_{\text{NCI at acquisition}}+58.8+2.2
      -\underbrace{31.2}_{\text{NCI at disposal}}=353.8.
\]
Parent equity also follows directly from \(300+53.8=353.8\). In the full
sale, the parent's separate cost-based investment gain is
\(169-115.2=53.8\), already equal to its cumulative consolidated profit
in this example. Adding that separate gain to the group result would count
the investment return twice.

Intermediate statements are essential checks. Freezing goodwill at USD 19.2
understates both goodwill and parent OCI by 1.6 before disposal. Yet its
larger pre-recycling disposal gain and smaller recycled reserve cancel,
leaving the same final gain of 33.8. Dividing goodwill translation 80/20
would misattribute 0.32 to NCI. Recycling all 12.6 while still removing NCI
31.2 would overstate profit by 2.2 and leave a spurious negative parent
reserve; recycling only the sold fraction in the retained-interest case
would leave 2.6 after control is lost. Final total equity alone misses these
different failures.

The case uses relative event dates and historical technical support. Actual
valuation, tax, impairment, acquisition exceptions and applicable amendments
remain required inputs. A retained 20\% interest also needs a separate
significant-influence and classification assessment before subsequent
measurement; its percentage does not supply that conclusion. The same
separation of legal-entity balances and group adjustments explains why
eliminating an intragroup loan can leave a currency exposure.
